\documentclass[11pt]{article}
\usepackage[margin=25mm]{geometry}
\usepackage{fontspec}
\setmainfont{Latin Modern Roman}
\usepackage{amsmath,amssymb,mathtools}
\usepackage{unicode-math}
\setmathfont{Latin Modern Math}
\usepackage{xurl,hyperref}
\hypersetup{colorlinks=true,urlcolor=blue}
\setcounter{secnumdepth}{0}
\setlength{\parindent}{0pt}
\setlength{\parskip}{5pt}
\setlength{\emergencystretch}{3em}
\newcommand{\sref}[2]{\hyperlink{source:#1:#2}{[S#2]}}
\newcommand{\cataloguescope}[1]{\paragraph{Recorded target.}#1}
\begin{document}
\section{Generalized Star Height Problem}\label{statement}

\subsection{Definitions and mathematical statement}

\paragraph{Shared notation.}$\mathbb N=\{1,2,\ldots\}$, and $\mathbb Z,\mathbb Q,\mathbb R,\mathbb C$ denote the integers, rationals, reals and complex numbers. Any other conventions are specified locally.


Let $\Sigma$ be a finite nonempty alphabet and $\Sigma^*$ its finite words, including the empty word $\varepsilon$. A generalized regular expression starts with $\varnothing$, $\varepsilon$ and letters $a\in\Sigma$, and combines expressions using union, concatenation, Kleene star and complement relative to $\Sigma^*$. For languages $L,M$, concatenation is $LM=\{uv:u\in L,v\in M\}$ and $L^*=\bigcup_{j\ge0}L^j$, with $L^0=\{\varepsilon\}$. Its star height is defined by
\[
 h(\varnothing)=h(\varepsilon)=h(a)=0,\quad
 h(E\cup F)=h(EF)=\max\{h(E),h(F)\},\quad
 h(E^*)=h(E)+1,\quad h(E^{c})=h(E).
\]
A language is regular if some such expression represents it, equivalently if it is recognized by a finite automaton. Its generalized star height is $h(L)=\min\{h(E):E\text{ represents }L\}$; complement is allowed in this minimum.
\paragraph{Statement recorded in the supplied source.}
Is there a single integer $B\ge0$ such that $h(L)\le B$ for every regular language $L$ over every finite alphabet? The bound must be independent of the language and alphabet. The frequently studied assertion $h(L)\le1$ for all regular languages is a possible stronger answer, but the catalogue's question asks whether any universal finite bound exists. \sref{statement}{2}
\subsection{Short English statement}
Is there a universal finite limit on nested Kleene stars needed to describe regular languages when complement is also allowed?
\subsection{Sources}
\begin{description}
\item[\hypertarget{source:statement:1}{S1}] ULAM.AI and the UnsolvedMath Contributors, UnsolvedMath: A Curated Collection of Open Mathematics Problems, catalogue entry. CC BY 4.0. \url{https://huggingface.co/datasets/ulamai/UnsolvedMath}.
\item[\hypertarget{source:statement:2}{S2}] Tom Bourne and Nik Ruškuc, On the Star-Height of Subword Counting Languages and Their Relationship to Rees Zero-Matrix Semigroups, arXiv:1603.06236v2 (2016), Section 1, definitions and the question of generalized height greater than one. \url{https://arxiv.org/pdf/1603.06236v2}.
\end{description}
\label{end:statement}
\end{document}
